\section{Training the TextCNN Reasoning-Quality Classifier}
\label{sec:textcnn-judge-training}

This document summarizes how we construct the GOOD/BAD reasoning-quality dataset
and train the TextCNN classifier used for ACRD ML filtering. It is based on
\texttt{minerva-judge/DATASET.md} and \texttt{minerva-judge/CLASSIFIER\_TRAINING.md}.

\subsection{Dataset Curation}
\label{subsec:textcnn-dataset-curation}

\paragraph{Source prompts.}
We start from the Minerva base train split:
\texttt{dataset/minerva\_base\_split/minerva-base-train.jsonl}.
Each row contributes the prompt, answer/ground\_truth, reward\_fn, and task, and
is assigned a stable UID (hash of prompt) if missing.

\paragraph{Response generation (plain + hinted).}
We generate model responses in two variants per prompt using
\texttt{minerva-judge/scripts/generate\_answer\_guided.py}:
\begin{itemize}
  \item \textbf{plain}: original prompt only.
  \item \textbf{hinted}: prompt plus an ACR block containing
        \texttt{GROUND\_TRUTH\_LABELS}, \texttt{LABEL\_REFERENCE}, and the marker
        text ``You are generating a reasoning trace for training.''.
\end{itemize}
The ACR block is built by \texttt{minerva.acr\_prompt.build\_acr\_block} with
label details supplied by \texttt{LabelDetailsStore}. Label references are
truncated to 4048 characters if needed; if the prompt becomes too long, the
reference is omitted. Unless \texttt{--no-system-prompt} is set, we include the
CTI system prompt from \texttt{rlvr/mydata/data\_prepare/prompt.py}.

Models used for response generation:
\begin{itemize}
  \item \texttt{meta-llama/Llama-3.1-8B-Instruct}
  \item \texttt{meta-llama/Llama-3.2-3B-Instruct}
  \item \texttt{Qwen/Qwen3-4B-Instruct-2507}
  \item \texttt{Qwen/Qwen3-8B}
  \item \texttt{openai/gpt-oss-20b}
\end{itemize}
Typical generation settings: temperature 0.7, max\_new\_tokens 1024, both variants.
Outputs are written to \texttt{minerva-judge/data/responses\_<model>.jsonl}.

\paragraph{Reward filtering (correct-only).}
Each response is scored with \texttt{reward\_minerva}. We define
\texttt{correct = (reward >= 0.999)} and keep only correct responses for
judge labeling and classifier training, isolating reasoning quality from answer
correctness.

\paragraph{LLM judge labeling.}
Correct responses are labeled by a rubric-based LLM judge
(\texttt{openai/gpt-oss-120b}) via
\texttt{minerva-judge/scripts/score\_with\_judge.py}. We do not provide a system
prompt by default. The \textbf{exact} judge prompt template is:

\begin{verbatim}
You are an expert CTI evaluator. You will be given a Question and a Response.
The Question may include answer/reference text as hints.

Label the Response as GOOD or BAD. Mark BAD if any single criterion below is present;
otherwise mark GOOD. Do not grade correctness beyond these checks.

Return exactly one JSON object on a single line and nothing else.
- GOOD: {"label": "GOOD"}
- BAD: {"label": "BAD", "category_id": <number>, "category_title": "<title>"}

If multiple BAD criteria apply, choose the lowest-numbered category.

Criteria (one per bullet):
1. Leakage: says or implies the answer/label/options/reference were provided,
   or quotes/paraphrases provided reference text instead of reasoning from the prompt.
2. Incoherent: loops, repeated phrases/lines, templated filler, or gibberish.
3. Ungrounded: invents concrete details not in the question (extra CVEs,
   vendors, malware names, IOCs, dates, techniques, etc.).
4. Mismatch: reasoning supports a different label than the final answer,
   or directly contradicts it.
5. Unsupported: provides reasoning but does not use evidence from the prompt to
   justify the answer (hand-wavy or irrelevant justification), OR provides an
   answer with no reasoning at all (answer-only).
6. Other: refusals, policy/meta artifacts, generic CTI tutorials, or prompt copying.

Example outputs:
{"label": "GOOD"}
{"label": "BAD", "category_id": 1, "category_title": "Leakage"}


QUESTION:
{QUESTION}

RESPONSE:
{RESPONSE}
\end{verbatim}

We exclude \texttt{judged\_gpt\_oss\_120b.jsonl} from classifier training to
avoid overlap with the judge model itself.

\paragraph{Synthetic BAD augmentation.}
To increase coverage of BAD categories, we synthesize additional BAD responses
using \texttt{openai/gpt-oss-120b}. Each synthetic example is conditioned on a
hinted prompt and a randomly sampled rubric violation criterion. We then re-score
these responses and keep only reward-correct synthetic rows.

\paragraph{Dataset assembly and split.}
The builder \texttt{minerva-judge/classifier/build\_dataset.py}:
\begin{enumerate}
  \item Joins judged rows back to responses using a stable key.
  \item Keeps only \texttt{judge\_label} in \{GOOD, BAD\} and \texttt{correct=true}.
  \item Adds synthetic BADs (correct-only).
  \item Balances via downsampling to 32k GOOD + 32k BAD (seed 1337).
  \item Splits 80/20 by class into train/val.
\end{enumerate}
Outputs:
\texttt{minerva-judge/classifier/data/train.jsonl} and
\texttt{minerva-judge/classifier/data/val.jsonl}.

\paragraph{Row joining and response selection details.}
Judge rows are joined to response rows by a stable key
\texttt{uid:model:variant[:attempt]}. For response text, the builder prefers
\texttt{response\_final} if present; otherwise it strips trailing markers (e.g.,
\texttt{assistantfinal}) from the raw response and falls back to the raw response.
This ensures the classifier sees clean response text.

\subsection{TextCNN Training}
\label{subsec:textcnn-training}

\paragraph{Model input and tokenization.}
TextCNN uses regex tokenization \texttt{[A-Za-z0-9\_]+} with lowercasing.
Default text mode is \texttt{response\_prompt}, which concatenates:
response tokens, a literal \texttt{sep}, then prompt tokens. If a max token
budget is enforced, response tokens are kept in full when possible and prompt
tokens are truncated to the remaining budget.

Vocabulary is built from the training split with \texttt{max\_vocab=100000}
and \texttt{min\_freq=2}. Special tokens are \texttt{<pad>} and \texttt{<unk>}.

\paragraph{Architecture.}
The TextCNN classifier is a standard 1D CNN over token embeddings:
\begin{itemize}
  \item Embedding layer (random init unless pretrained vectors are provided).
  \item 1D convolutions with kernel sizes (default 3,4,5), ReLU, max-over-time pooling.
  \item Concatenation, dropout, and a 2-way linear classifier.
\end{itemize}

\paragraph{Training setup.}
Training uses AdamW with cross-entropy loss, batch size 128, and 5 epochs. We
report validation accuracy/precision/recall/F1 after each epoch; the final
metrics are written to \texttt{metrics.json} in the output directory.

\paragraph{Hyperparameter sweep.}
We sweep the following grid (run via
\texttt{minerva-judge/classifier/scripts/run\_all\_textcnn\_parallel.sh}):
\begin{itemize}
  \item Learning rate: \{3e-4, 6e-4, 1e-3, 2e-3, 4e-3\}
  \item Kernel sizes: \{3,4,5\}
  \item Filters per kernel: \{256, 384\}
  \item Dropout: \{0, 0.25\}
  \item Embedding dim: \{200, 300\}
  \item Max tokens: \{2048, 3072\}
\end{itemize}

\paragraph{Best model (used for ACRD ML filtering).}
The best TextCNN model reported in the sweep is:
\begin{itemize}
  \item \texttt{textcnn\_lr6e-4\_k345\_f384\_e200\_t3072\_d0p25}
  \item lr=6e-4, kernels=3/4/5, filters=384, embed\_dim=200,
        max\_tokens=3072, dropout=0.25
  \item Validation accuracy: 0.828125
\end{itemize}

\subsection{Notes}
\label{subsec:textcnn-notes}

\begin{itemize}
  \item Classifier labels are derived only from reward-correct responses.
  \item Synthetic BADs are included only if their extracted answers are reward-correct.
\end{itemize}
